% =====================================================================
% core_model.tex — replaces Section 4 ("Structural model") of the paper.
%
% % =====================================================================
% core_model.tex — replaces Section 4 ("Structural model") of the paper.
%
% % =====================================================================
% core_model.tex — replaces Section 4 ("Structural model") of the paper.
%
% % =====================================================================
% core_model.tex — replaces Section 4 ("Structural model") of the paper.
%
% \input{core_model} from the main file. Requires amsmath, amssymb and the
% \citep/\citet macros already used in the paper. Cross-references assume the
% appendix file model_appendix.tex is also included (labels app:model-*).
%
% What changed relative to the previous version of this section:
%   (i)  upstream sectors now contain a FINITE number N_s of varieties, so the
%        sourcing extensive margin is a genuine random object and observed
%        zeros are model predictions rather than dropped cells;
%   (ii) comparative advantage is located at the ATTRACTION-AREA level, the
%        level at which the extensive-margin regression identifies the
%        trade-cost elasticity;
%  (iii) trade costs are the power law tau = d^alpha, so the extensive-margin
%        distance coefficient maps into alpha with no attenuation correction.
% =====================================================================

\section{Structural model}
\label{sec:model}

We next develop a structural model that rationalizes the spatial propagation patterns
documented in the data and allows us to quantify the role of production networks in shaping
the incidence of sectoral shocks. The model has a Ricardian structure that is used to recover
firm-to-firm linkages not directly observed in the data. Based on the predicted production
network, we then simulate shocks to the downstream industry and quantify their spatial
diffusion.

Relative to a textbook Eaton--Kortum economy the model departs on one dimension: upstream
sectors contain a \emph{finite} number of varieties. This is what makes the sourcing extensive
margin --- the object our survey data measure --- a genuine random outcome. With a continuum of
varieties every region with a positive comparative advantage supplies every buyer, so the
zeros that dominate our data could only be accommodated by discarding the corresponding cells.
Here they are model predictions.

\subsection{Model}
\label{sec:model-model}

\paragraph{Environment.} The economy consists of $R+1$ regions, where the $(R+1)$th region
represents the rest of the world. Firms operate either in upstream sectors $s=1,\dots,S$ or in
the downstream industry $d$ (motor vehicles or aerospace). Each domestic region hosts at most
one representative downstream firm. Firm locations and sectoral assignments are taken as
exogenous. We allow for imports, but in the absence of balance-sheet information on foreign
suppliers we do not model firm-level linkages abroad; the foreign share of sourcing is
calibrated (Appendix~\ref{app:model-foreign}).

Domestic regions are additionally grouped into \emph{attraction areas}. An attraction area $a$
is the set of upstream regions sharing the same nearest downstream producer, and we write
$a(r')$ for the area of region $r'$ and $\mathcal{R}_{a}$ for its members. This is the same
partition as the fixed effect of the extensive-margin regression in Equation~(1), and that is
not a coincidence: it is the level at which the data identify the trade-cost elasticity.

Because the model is estimated separately for the motor vehicles and aerospace industries, we
omit the downstream-industry index in the remainder of the exposition. Unless otherwise stated,
all variables and parameters should therefore be understood as industry-specific.

\paragraph{Final demand.} Downstream firms compete under monopolistic competition and face an
isoelastic demand function:
\begin{equation}
q_{r}=\left(\frac{p_{r}}{P}\right)^{\varepsilon-1}\frac{E}{P}\,\delta^{D}_{r},
\end{equation}
where $\delta^{D}_{r}$ is a region-specific demand shifter, $E/P$ denotes aggregate real
expenditure and $p_{r}/P$ is the relative price of the variety produced in region $r$. Demand
shifters are later used to simulate downstream shocks and trace their propagation through
production networks. The market for final goods is assumed to be perfectly integrated.

\paragraph{Production.} Downstream firms produce using labor and intermediate inputs according
to a constant-returns-to-scale production technology:
\begin{equation}
y_{r}=A_{r}\,g\big(\ell_{r},\{q_{rs}\}_{s=1}^{S}\big),
\label{eq:downstream-tech}
\end{equation}
where $A_{r}$ denotes total factor productivity, $\ell_{r}$ is labor input and $q_{rs}$ is the
quantity of the input bundle from sector $s$ used in production. As discussed in
Appendix~\ref{app:model-derivation}, the structural estimation assumes a nested CES
specification, where $\Omega_{L}$ scales the labor share, $\Omega_{s}$ denotes the expenditure
share on inputs from sector $s$, and each $q_{rs}$ is itself a CES aggregate over the varieties
of sector $s$.

\paragraph{Upstream technology.} Each upstream sector $s$ consists of $N_{s}$ horizontally
differentiated varieties, indexed $\rho=1,\dots,N_{s}$. Within each variety and each region
there is a continuum of firms whose productivities above $z$ form a Poisson process with mean
$T_{r's}z^{-\theta}$. The most productive firm of region $r'$ for variety $\rho$ --- its
\emph{champion} --- therefore satisfies
\begin{equation}
\Pr\big(\varphi_{r'\rho s}\le z\big)=e^{-T_{r's}z^{-\theta}},
\qquad T_{r's}>0,
\label{eq:frechet}
\end{equation}
so one Fr\'echet draw per (region, variety) pair \emph{is} the champion of a continuum of
firms rather than an approximation to it. Because the iceberg cost scales all firms of a region
identically, that champion is the only firm of region $r'$ that ever serves a buyer sourcing
$\rho$ from $r'$: a firm in the model is a (region, variety) champion. As in Eaton and Kortum
(2002), $T_{r's}$ governs average productivity while $\theta$ controls the dispersion of
productivity draws and hence the strength of comparative advantage. Only the number of
\emph{varieties} is finite; the firms are not discretized.

\paragraph{Comparative advantage at the area level.} Comparative advantage is a property of the
attraction area rather than of the individual region:
\begin{equation}
T_{r's}=T_{a(r')s}\qquad\text{for every }r'\in\mathcal{R}_{a}.
\label{eq:T-area}
\end{equation}
Productivity draws remain independent across regions: two regions in the same area share the
Fr\'echet \emph{scale}, not the realization of their champions, so each still runs its own
Ricardian competition for every variety. What \eqref{eq:T-area} removes is unobserved
technological heterogeneity \emph{within} an area, and the payoff is that within an area the
only thing separating a region that hosts a supplier from one that does not is distance. This
is what makes the model explain the same variation the regression uses.

\paragraph{Sourcing.} Shipping intermediate goods from region $r'$ to region $r$ involves
iceberg trade costs $\tau_{r'rs}$, which we parameterize as a power law in geographic distance,
\begin{equation}
\tau_{r'rs}=d_{r'r}^{\,\alpha}.
\label{eq:tau}
\end{equation}
Following Eaton and Kortum (2002), downstream firms source each variety from the lowest-cost
supplier, in France or abroad. The probability that a downstream firm located in region $r$
sources a given variety of sector $s$ from region $r'$ is
\begin{equation}
\gamma_{r'rs}=\frac{T_{a(r')s}\,(w_{r's}\tau_{r'rs})^{-\theta}}
{\sum_{k=1}^{R+1}T_{a(k)s}\,(w_{ks}\tau_{krs})^{-\theta}}
\;\equiv\;\frac{T_{a(r')s}\,(w_{r's}\tau_{r'rs})^{-\theta}}{\Phi_{rs}},
\label{eq:ek-share}
\end{equation}
where $w_{r's}$ denotes the unit production cost in region--sector pair $(r',s)$. Because trade
costs increase with geographic distance, sourcing probabilities decline with spatial
separation, generating localized production networks. Two regions $r',r''$ in the same area
facing the same buyer satisfy
\begin{equation}
\ln\frac{\gamma_{r'rs}}{\gamma_{r''rs}}
=-\theta\alpha\big(\ln d_{r'r}-\ln d_{r''r}\big)-\theta\ln\frac{w_{r's}}{w_{r''s}},
\label{eq:within-area}
\end{equation}
the technology term and the price index having cancelled: the within-area sourcing profile is a
function of $(\theta,\alpha)$ and observed distances alone.

\paragraph{The extensive margin.} Our surveys record whether an upstream firm supplies the
downstream industry, not which buyer it serves. The relevant object is therefore not the
bilateral sourcing probability \eqref{eq:ek-share} but the probability that a region wins a
given variety \emph{somewhere} in the industry --- the union over downstream buyers,
\begin{equation}
q_{r's}\;\equiv\;\Pr\Big(\bigcup_{r=1}^{R}\big\{\text{$r'$ wins variety $\rho$ for buyer $r$}\big\}\Big).
\label{eq:q}
\end{equation}
With $N_{s}$ varieties the win events are independent across $\rho$,\footnote{Independence
across varieties follows from the granularity assumption alone (Appendix~\ref{app:model-derivation}):
the champion draws
$\varphi_{r'\rho s}$ are independent across $\rho$, and which origin wins a variety is a
comparison of \emph{relative} delivered costs, so the winner of one variety carries no
information about the winner of another. Note what this does \emph{not} require. It is silent
on the price index, because the identity of the winner is settled before any aggregation: the
sourcing decision is a per-variety cost comparison, whereas $p_{rs}$ is an aggregator built
once the winners are known. What the price index does affect is the \emph{value} attached to
each win, since expenditure shares depend on $p_{rs}(\rho)/p_{rs}$ and the realized index over
$N_{s}$ varieties is itself random. There we follow \citet{lmm2023} and let downstream firms
optimize against the expected (certainty-equivalent) index rather than internalize how the
selection of any one supplier moves it, which keeps the success probability common and
non-random across varieties in the value block as well. Appendix~\ref{app:model-price}
quantifies what this assumption costs and why it does not touch the identification of $\alpha$
or $N_{s}$.} so the industry-wide supplier count is binomial,
\begin{equation}
K_{r's}\sim\mathrm{Bin}\big(N_{s},q_{r's}\big),
\qquad
\Pr\big(K_{r's}=0\big)=(1-q_{r's})^{N_{s}}>0 .
\label{eq:binomial}
\end{equation}
The zero-supplier probability is strictly positive and endogenous: it falls with comparative
advantage, rises with distance through \eqref{eq:tau}, and falls with the variety count. Under
a continuum it is zero for every region with $T_{r's}>0$, which is the sense in which the
extensive margin previously had no content.

The union in \eqref{eq:q} is not the product of its per-buyer marginals. For a given variety
the champion cost is common across buyers and only $\tau_{r'rs}$ differs, so win events are
\emph{positively} dependent across buyers and $q_{r's}>1-\prod_{r}(1-\gamma_{r'rs})$: a firm
productive enough to win one market is more likely to win the others. The per-buyer marginal
$K_{r'rs}\sim\mathrm{Bin}(N_{s},\gamma_{r'rs})$ remains exact, but $q_{r's}$ has no tractable
closed form, and we evaluate it by simulation. This is why the model is estimated by simulated
rather than analytical method of moments.\footnote{Writing $q_{r's}$ in closed form requires
inclusion--exclusion over subsets of downstream buyers,
$q_{r's}=\sum_{\emptyset\neq\mathcal{S}\subseteq\{1,\dots,R\}}(-1)^{|\mathcal{S}|+1}
\Pr(\cap_{r\in\mathcal{S}}\{\text{$r'$ wins for }r\})$, a sum with $2^{R}-1$ terms. With the
number of downstream regions we observe this is prohibitive, and dominance pruning of the
subsets --- discarding buyers whose win event contains another's --- reduces the count but not
reliably enough to be usable inside an optimizer. The alternative of a closed form is to impose
independence and use $1-\prod_{r}(1-\gamma_{r'rs})$, which is a strict lower bound on
\eqref{eq:q} and biases the extensive margin; simulation evaluates the union exactly, at the
cost of simulation noise that the weighting matrix accounts for.}

\paragraph{Intuition: what the model implies for the extensive-margin regression.} The union
\eqref{eq:q} has no closed form, but the single-buyer case does, and it is what gives the
empirical specification of Equation~(1) its structural reading. Consider a firm of productivity
$z$ located in $r'$ facing only its area's anchor buyer $a$. Its probability of \emph{not}
supplying is complementary log-log,
\begin{equation}
\Pr\big(\text{no supply}\mid z\big)=1-\exp\big(-\exp(\eta)\big),
\qquad
\eta=\underbrace{\ln\Phi^{-r'}_{as}}_{\text{area}\times\text{sector FE}}
+\,\theta\ln w_{r's}
+\,\theta\alpha\,\ln d_{r'a}
\;-\;\theta\,\ln z ,
\label{eq:cloglog}
\end{equation}
where $\Phi^{-r'}_{as}$ collects the competitors' cost scale (Appendix~\ref{app:model-prop}
derives \eqref{eq:cloglog} and states the conditions under which it is exact). Three features
carry over to the estimation. The distance coefficient is $+\theta\alpha$, which is what makes
the extensive-margin regression informative about the trade-cost elasticity. The variety count
is \emph{absent}: conditioning on the firm removes it, so $\alpha$ and $N_{s}$ load on separate
features of the data. And the coefficient on own productivity is $-\theta$, a free
over-identifying test of a parameter we otherwise calibrate.

With several buyers the outcome is the union and \eqref{eq:cloglog} is no longer exact. We
therefore do not read $\alpha$ off the empirical coefficient directly. Instead we treat the
regression as an auxiliary model: we run the \emph{same} cloglog on simulated and observed
data, where the simulated outcome is the union event \eqref{eq:q}, and match the coefficients.
This is consistent for $\alpha$ whether or not the link is exact, and it is why the distance
coefficient enters the moment vector rather than being inverted analytically.

\subsection{Estimation and fit}
\label{sec:model-estimation}

We estimate the model using Simulated Method of Moments. Unlike the standard Ricardian trade
literature, which relies on bilateral input--output flows, such flows are not observed in our
setting. Instead, we exploit survey information on supplier relationships together with VAT
data on firm sales to discipline the spatial structure of sourcing and demand.

The model's parameters are
\begin{equation}
\big\{\Omega_{L},\;\{\Omega_{s}\},\;\alpha,\;\{A_{r}\}\big\}
\;\cup\;\big\{\{T_{as}\},\;\{N_{s}\}\big\},
\end{equation}
but only the first group is searched over. Comparative advantage and the variety count are
\emph{concentrated out}: given any candidate value of $(\Omega_{L},\{\Omega_{s}\},\alpha,\{A_{r}\})$,
each is pinned by the model itself rather than left free.

This is not a computational shortcut but a property of the moment structure. For $\{T_{as}\}$,
the area-aggregated sourcing shares \eqref{eq:gamma-area} are as numerous as the comparative-advantage
parameters, and for given trade costs and expenditures the mapping from $T$ to those shares is
invertible: there is a unique vector of comparative advantages, up to the per-sector reference
normalization, that reproduces the observed shares exactly. We compute it by a matrix-scaling
iteration alternated with the general-equilibrium solve
(Appendix~\ref{app:model-estimation}). For $\{N_{s}\}$, the count moment $\bar G_{s}(0)$ is
decreasing in the variety count and, by the argument below, does not require re-simulation, so
the variety count that best matches it is found by a monotone integer search over the bounds
\eqref{eq:app-bounds}. Both are exact concentrations in the sense of a profiled extremum
estimator: minimizing the criterion jointly over all parameters and minimizing it over
$(\Omega_{L},\{\Omega_{s}\},\alpha,\{A_{r}\})$ with $T$ and $N$ set to their concentrated values
give the same estimate.

Two things follow. The dimension of the outer search collapses from several hundred parameters
to a handful, which is what makes the estimation feasible. And the division of labour in
identification becomes explicit: the sourcing-share block is exhausted by $T$ and the count
block by $N$, so the extensive-margin coefficient is fit by $\alpha$ alone. The over-identifying
content of the area-level restriction is preserved, since a single $T_{as}$ must rationalize the
shares of every region in the area.

Beyond these, the model includes several demand- and supply-side elasticities, which we
calibrate as described in Appendix~\ref{app:model-calibration}. We also provide there
additional details on the optimization procedure and on the calculation of standard errors.

\paragraph{Targeted moments.} The model is matched to six families of moments. The first three
are as before. The downstream industry's labor share and the sectoral composition of
intermediate expenditures are constructed from input-output tables and identify $\Omega_{L}$
and $\{\Omega_{s}\}_{s=1}^{S}$. The regional distribution of downstream sales, $\pi_{r}$, is
recovered from the VAT data and identifies the downstream productivity parameters $\{A_{r}\}$,
up to a multiplicative constant and therefore normalized relative to the largest commuting
zone. The extensive-margin distance coefficient of Equation~(1) identifies the trade-cost
elasticity $\alpha$ through \eqref{eq:cloglog}.

The remaining two families are where the granular structure enters. \emph{Sourcing shares} are
now aggregated to the attraction area,
\begin{equation}
\tilde\gamma_{as}=\sum_{r'\in\mathcal{R}_{a}}\tilde\gamma_{r's},
\qquad
\tilde\gamma_{r's}=\sum_{r}\frac{X_{rs}}{\sum_{r}X_{rs}}\,\gamma_{r'rs},
\label{eq:gamma-area}
\end{equation}
where $X_{rs}$ denotes the value of intermediate inputs from sector $s$ shipped to downstream
firms located in region $r$, and identify the area-level comparative advantage $\{T_{as}\}$.
For each sector these are normalized relative to the area with the largest sourcing share.
\emph{Count moments} discipline the variety count: writing $\mathcal{R}^{+}_{s}$ for the set of
regions lying inside attraction areas that host at least one supplier of sector $s$, the share
of those regions hosting no supplier is
\begin{equation}
\bar G_{s}(0)=\frac{1}{|\mathcal{R}^{+}_{s}|}\,
\#\big\{r'\in\mathcal{R}^{+}_{s}:K_{r's}=0\big\},
\qquad
\mathbb{E}\big[\bar G_{s}(0)\big]
=\frac{1}{|\mathcal{R}^{+}_{s}|}\sum_{r'\in\mathcal{R}^{+}_{s}}(1-q_{r's})^{N_{s}} .
\label{eq:count-moment}
\end{equation}
This is the count analogue of the number-of-buyers-per-seller moment in \citet{lmm2023}. It is
decreasing in $N_{s}$, which is what identifies the variety count: given the sourcing shares
--- pinned by the area aggregates and by $\alpha$, neither of which involves $N_{s}$ --- the
only free object shaping \eqref{eq:count-moment} is the number of varieties. The support is
restricted to active areas because regions in areas with no supplier at all are structural
zeros, absorbed by the fixed effect of Equation~(1) and carrying no information about
$(\alpha,N_{s})$; Appendix~\ref{app:model-zeros} makes the distinction precise.

Two observations discipline the exercise further. The variety count is bounded by the observed
number of distinct suppliers: since each variety generates one supplying firm per winning
origin, $N^{\mathrm{obs}}_{s}/R^{D}\le N_{s}\le N^{\mathrm{obs}}_{s}$ where $R^{D}$ is the
number of downstream buyers, and matching the observed firm count gives a second, independent
estimate that must fall inside those bounds (Appendix~\ref{app:model-bounds}). And because
$T_{as}$ is constant within an area, the area aggregate \eqref{eq:gamma-area} no longer
exhausts the region-level shares: each active region beyond the first in an area becomes an
over-identifying restriction, requiring that the split of an area's sourcing across its regions
follow the distance profile \eqref{eq:within-area}.

\paragraph{Identification.} The channels are close to orthogonal by construction. $\alpha$ is
read from a \emph{within}-area gradient that comparative advantage cannot touch, because
comparative advantage is constant within the area. $T_{as}$ is read from \emph{across}-area
value shares. $N_{s}$ is read from the \emph{level} of the empty-region share, which the fixed
effect strips out of the $\alpha$ regression. And $\theta$, which we calibrate, is over-identified
by the own-productivity coefficient in \eqref{eq:cloglog}.
 from the main file. Requires amsmath, amssymb and the
% \citep/\citet macros already used in the paper. Cross-references assume the
% appendix file model_appendix.tex is also included (labels app:model-*).
%
% What changed relative to the previous version of this section:
%   (i)  upstream sectors now contain a FINITE number N_s of varieties, so the
%        sourcing extensive margin is a genuine random object and observed
%        zeros are model predictions rather than dropped cells;
%   (ii) comparative advantage is located at the ATTRACTION-AREA level, the
%        level at which the extensive-margin regression identifies the
%        trade-cost elasticity;
%  (iii) trade costs are the power law tau = d^alpha, so the extensive-margin
%        distance coefficient maps into alpha with no attenuation correction.
% =====================================================================

\section{Structural model}
\label{sec:model}

We next develop a structural model that rationalizes the spatial propagation patterns
documented in the data and allows us to quantify the role of production networks in shaping
the incidence of sectoral shocks. The model has a Ricardian structure that is used to recover
firm-to-firm linkages not directly observed in the data. Based on the predicted production
network, we then simulate shocks to the downstream industry and quantify their spatial
diffusion.

Relative to a textbook Eaton--Kortum economy the model departs on one dimension: upstream
sectors contain a \emph{finite} number of varieties. This is what makes the sourcing extensive
margin --- the object our survey data measure --- a genuine random outcome. With a continuum of
varieties every region with a positive comparative advantage supplies every buyer, so the
zeros that dominate our data could only be accommodated by discarding the corresponding cells.
Here they are model predictions.

\subsection{Model}
\label{sec:model-model}

\paragraph{Environment.} The economy consists of $R+1$ regions, where the $(R+1)$th region
represents the rest of the world. Firms operate either in upstream sectors $s=1,\dots,S$ or in
the downstream industry $d$ (motor vehicles or aerospace). Each domestic region hosts at most
one representative downstream firm. Firm locations and sectoral assignments are taken as
exogenous. We allow for imports, but in the absence of balance-sheet information on foreign
suppliers we do not model firm-level linkages abroad; the foreign share of sourcing is
calibrated (Appendix~\ref{app:model-foreign}).

Domestic regions are additionally grouped into \emph{attraction areas}. An attraction area $a$
is the set of upstream regions sharing the same nearest downstream producer, and we write
$a(r')$ for the area of region $r'$ and $\mathcal{R}_{a}$ for its members. This is the same
partition as the fixed effect of the extensive-margin regression in Equation~(1), and that is
not a coincidence: it is the level at which the data identify the trade-cost elasticity.

Because the model is estimated separately for the motor vehicles and aerospace industries, we
omit the downstream-industry index in the remainder of the exposition. Unless otherwise stated,
all variables and parameters should therefore be understood as industry-specific.

\paragraph{Final demand.} Downstream firms compete under monopolistic competition and face an
isoelastic demand function:
\begin{equation}
q_{r}=\left(\frac{p_{r}}{P}\right)^{\varepsilon-1}\frac{E}{P}\,\delta^{D}_{r},
\end{equation}
where $\delta^{D}_{r}$ is a region-specific demand shifter, $E/P$ denotes aggregate real
expenditure and $p_{r}/P$ is the relative price of the variety produced in region $r$. Demand
shifters are later used to simulate downstream shocks and trace their propagation through
production networks. The market for final goods is assumed to be perfectly integrated.

\paragraph{Production.} Downstream firms produce using labor and intermediate inputs according
to a constant-returns-to-scale production technology:
\begin{equation}
y_{r}=A_{r}\,g\big(\ell_{r},\{q_{rs}\}_{s=1}^{S}\big),
\label{eq:downstream-tech}
\end{equation}
where $A_{r}$ denotes total factor productivity, $\ell_{r}$ is labor input and $q_{rs}$ is the
quantity of the input bundle from sector $s$ used in production. As discussed in
Appendix~\ref{app:model-derivation}, the structural estimation assumes a nested CES
specification, where $\Omega_{L}$ scales the labor share, $\Omega_{s}$ denotes the expenditure
share on inputs from sector $s$, and each $q_{rs}$ is itself a CES aggregate over the varieties
of sector $s$.

\paragraph{Upstream technology.} Each upstream sector $s$ consists of $N_{s}$ horizontally
differentiated varieties, indexed $\rho=1,\dots,N_{s}$. Within each variety and each region
there is a continuum of firms whose productivities above $z$ form a Poisson process with mean
$T_{r's}z^{-\theta}$. The most productive firm of region $r'$ for variety $\rho$ --- its
\emph{champion} --- therefore satisfies
\begin{equation}
\Pr\big(\varphi_{r'\rho s}\le z\big)=e^{-T_{r's}z^{-\theta}},
\qquad T_{r's}>0,
\label{eq:frechet}
\end{equation}
so one Fr\'echet draw per (region, variety) pair \emph{is} the champion of a continuum of
firms rather than an approximation to it. Because the iceberg cost scales all firms of a region
identically, that champion is the only firm of region $r'$ that ever serves a buyer sourcing
$\rho$ from $r'$: a firm in the model is a (region, variety) champion. As in Eaton and Kortum
(2002), $T_{r's}$ governs average productivity while $\theta$ controls the dispersion of
productivity draws and hence the strength of comparative advantage. Only the number of
\emph{varieties} is finite; the firms are not discretized.

\paragraph{Comparative advantage at the area level.} Comparative advantage is a property of the
attraction area rather than of the individual region:
\begin{equation}
T_{r's}=T_{a(r')s}\qquad\text{for every }r'\in\mathcal{R}_{a}.
\label{eq:T-area}
\end{equation}
Productivity draws remain independent across regions: two regions in the same area share the
Fr\'echet \emph{scale}, not the realization of their champions, so each still runs its own
Ricardian competition for every variety. What \eqref{eq:T-area} removes is unobserved
technological heterogeneity \emph{within} an area, and the payoff is that within an area the
only thing separating a region that hosts a supplier from one that does not is distance. This
is what makes the model explain the same variation the regression uses.

\paragraph{Sourcing.} Shipping intermediate goods from region $r'$ to region $r$ involves
iceberg trade costs $\tau_{r'rs}$, which we parameterize as a power law in geographic distance,
\begin{equation}
\tau_{r'rs}=d_{r'r}^{\,\alpha}.
\label{eq:tau}
\end{equation}
Following Eaton and Kortum (2002), downstream firms source each variety from the lowest-cost
supplier, in France or abroad. The probability that a downstream firm located in region $r$
sources a given variety of sector $s$ from region $r'$ is
\begin{equation}
\gamma_{r'rs}=\frac{T_{a(r')s}\,(w_{r's}\tau_{r'rs})^{-\theta}}
{\sum_{k=1}^{R+1}T_{a(k)s}\,(w_{ks}\tau_{krs})^{-\theta}}
\;\equiv\;\frac{T_{a(r')s}\,(w_{r's}\tau_{r'rs})^{-\theta}}{\Phi_{rs}},
\label{eq:ek-share}
\end{equation}
where $w_{r's}$ denotes the unit production cost in region--sector pair $(r',s)$. Because trade
costs increase with geographic distance, sourcing probabilities decline with spatial
separation, generating localized production networks. Two regions $r',r''$ in the same area
facing the same buyer satisfy
\begin{equation}
\ln\frac{\gamma_{r'rs}}{\gamma_{r''rs}}
=-\theta\alpha\big(\ln d_{r'r}-\ln d_{r''r}\big)-\theta\ln\frac{w_{r's}}{w_{r''s}},
\label{eq:within-area}
\end{equation}
the technology term and the price index having cancelled: the within-area sourcing profile is a
function of $(\theta,\alpha)$ and observed distances alone.

\paragraph{The extensive margin.} Our surveys record whether an upstream firm supplies the
downstream industry, not which buyer it serves. The relevant object is therefore not the
bilateral sourcing probability \eqref{eq:ek-share} but the probability that a region wins a
given variety \emph{somewhere} in the industry --- the union over downstream buyers,
\begin{equation}
q_{r's}\;\equiv\;\Pr\Big(\bigcup_{r=1}^{R}\big\{\text{$r'$ wins variety $\rho$ for buyer $r$}\big\}\Big).
\label{eq:q}
\end{equation}
With $N_{s}$ varieties the win events are independent across $\rho$,\footnote{Independence
across varieties follows from the granularity assumption alone (Appendix~\ref{app:model-derivation}):
the champion draws
$\varphi_{r'\rho s}$ are independent across $\rho$, and which origin wins a variety is a
comparison of \emph{relative} delivered costs, so the winner of one variety carries no
information about the winner of another. Note what this does \emph{not} require. It is silent
on the price index, because the identity of the winner is settled before any aggregation: the
sourcing decision is a per-variety cost comparison, whereas $p_{rs}$ is an aggregator built
once the winners are known. What the price index does affect is the \emph{value} attached to
each win, since expenditure shares depend on $p_{rs}(\rho)/p_{rs}$ and the realized index over
$N_{s}$ varieties is itself random. There we follow \citet{lmm2023} and let downstream firms
optimize against the expected (certainty-equivalent) index rather than internalize how the
selection of any one supplier moves it, which keeps the success probability common and
non-random across varieties in the value block as well. Appendix~\ref{app:model-price}
quantifies what this assumption costs and why it does not touch the identification of $\alpha$
or $N_{s}$.} so the industry-wide supplier count is binomial,
\begin{equation}
K_{r's}\sim\mathrm{Bin}\big(N_{s},q_{r's}\big),
\qquad
\Pr\big(K_{r's}=0\big)=(1-q_{r's})^{N_{s}}>0 .
\label{eq:binomial}
\end{equation}
The zero-supplier probability is strictly positive and endogenous: it falls with comparative
advantage, rises with distance through \eqref{eq:tau}, and falls with the variety count. Under
a continuum it is zero for every region with $T_{r's}>0$, which is the sense in which the
extensive margin previously had no content.

The union in \eqref{eq:q} is not the product of its per-buyer marginals. For a given variety
the champion cost is common across buyers and only $\tau_{r'rs}$ differs, so win events are
\emph{positively} dependent across buyers and $q_{r's}>1-\prod_{r}(1-\gamma_{r'rs})$: a firm
productive enough to win one market is more likely to win the others. The per-buyer marginal
$K_{r'rs}\sim\mathrm{Bin}(N_{s},\gamma_{r'rs})$ remains exact, but $q_{r's}$ has no tractable
closed form, and we evaluate it by simulation. This is why the model is estimated by simulated
rather than analytical method of moments.\footnote{Writing $q_{r's}$ in closed form requires
inclusion--exclusion over subsets of downstream buyers,
$q_{r's}=\sum_{\emptyset\neq\mathcal{S}\subseteq\{1,\dots,R\}}(-1)^{|\mathcal{S}|+1}
\Pr(\cap_{r\in\mathcal{S}}\{\text{$r'$ wins for }r\})$, a sum with $2^{R}-1$ terms. With the
number of downstream regions we observe this is prohibitive, and dominance pruning of the
subsets --- discarding buyers whose win event contains another's --- reduces the count but not
reliably enough to be usable inside an optimizer. The alternative of a closed form is to impose
independence and use $1-\prod_{r}(1-\gamma_{r'rs})$, which is a strict lower bound on
\eqref{eq:q} and biases the extensive margin; simulation evaluates the union exactly, at the
cost of simulation noise that the weighting matrix accounts for.}

\paragraph{Intuition: what the model implies for the extensive-margin regression.} The union
\eqref{eq:q} has no closed form, but the single-buyer case does, and it is what gives the
empirical specification of Equation~(1) its structural reading. Consider a firm of productivity
$z$ located in $r'$ facing only its area's anchor buyer $a$. Its probability of \emph{not}
supplying is complementary log-log,
\begin{equation}
\Pr\big(\text{no supply}\mid z\big)=1-\exp\big(-\exp(\eta)\big),
\qquad
\eta=\underbrace{\ln\Phi^{-r'}_{as}}_{\text{area}\times\text{sector FE}}
+\,\theta\ln w_{r's}
+\,\theta\alpha\,\ln d_{r'a}
\;-\;\theta\,\ln z ,
\label{eq:cloglog}
\end{equation}
where $\Phi^{-r'}_{as}$ collects the competitors' cost scale (Appendix~\ref{app:model-prop}
derives \eqref{eq:cloglog} and states the conditions under which it is exact). Three features
carry over to the estimation. The distance coefficient is $+\theta\alpha$, which is what makes
the extensive-margin regression informative about the trade-cost elasticity. The variety count
is \emph{absent}: conditioning on the firm removes it, so $\alpha$ and $N_{s}$ load on separate
features of the data. And the coefficient on own productivity is $-\theta$, a free
over-identifying test of a parameter we otherwise calibrate.

With several buyers the outcome is the union and \eqref{eq:cloglog} is no longer exact. We
therefore do not read $\alpha$ off the empirical coefficient directly. Instead we treat the
regression as an auxiliary model: we run the \emph{same} cloglog on simulated and observed
data, where the simulated outcome is the union event \eqref{eq:q}, and match the coefficients.
This is consistent for $\alpha$ whether or not the link is exact, and it is why the distance
coefficient enters the moment vector rather than being inverted analytically.

\subsection{Estimation and fit}
\label{sec:model-estimation}

We estimate the model using Simulated Method of Moments. Unlike the standard Ricardian trade
literature, which relies on bilateral input--output flows, such flows are not observed in our
setting. Instead, we exploit survey information on supplier relationships together with VAT
data on firm sales to discipline the spatial structure of sourcing and demand.

The model's parameters are
\begin{equation}
\big\{\Omega_{L},\;\{\Omega_{s}\},\;\alpha,\;\{A_{r}\}\big\}
\;\cup\;\big\{\{T_{as}\},\;\{N_{s}\}\big\},
\end{equation}
but only the first group is searched over. Comparative advantage and the variety count are
\emph{concentrated out}: given any candidate value of $(\Omega_{L},\{\Omega_{s}\},\alpha,\{A_{r}\})$,
each is pinned by the model itself rather than left free.

This is not a computational shortcut but a property of the moment structure. For $\{T_{as}\}$,
the area-aggregated sourcing shares \eqref{eq:gamma-area} are as numerous as the comparative-advantage
parameters, and for given trade costs and expenditures the mapping from $T$ to those shares is
invertible: there is a unique vector of comparative advantages, up to the per-sector reference
normalization, that reproduces the observed shares exactly. We compute it by a matrix-scaling
iteration alternated with the general-equilibrium solve
(Appendix~\ref{app:model-estimation}). For $\{N_{s}\}$, the count moment $\bar G_{s}(0)$ is
decreasing in the variety count and, by the argument below, does not require re-simulation, so
the variety count that best matches it is found by a monotone integer search over the bounds
\eqref{eq:app-bounds}. Both are exact concentrations in the sense of a profiled extremum
estimator: minimizing the criterion jointly over all parameters and minimizing it over
$(\Omega_{L},\{\Omega_{s}\},\alpha,\{A_{r}\})$ with $T$ and $N$ set to their concentrated values
give the same estimate.

Two things follow. The dimension of the outer search collapses from several hundred parameters
to a handful, which is what makes the estimation feasible. And the division of labour in
identification becomes explicit: the sourcing-share block is exhausted by $T$ and the count
block by $N$, so the extensive-margin coefficient is fit by $\alpha$ alone. The over-identifying
content of the area-level restriction is preserved, since a single $T_{as}$ must rationalize the
shares of every region in the area.

Beyond these, the model includes several demand- and supply-side elasticities, which we
calibrate as described in Appendix~\ref{app:model-calibration}. We also provide there
additional details on the optimization procedure and on the calculation of standard errors.

\paragraph{Targeted moments.} The model is matched to six families of moments. The first three
are as before. The downstream industry's labor share and the sectoral composition of
intermediate expenditures are constructed from input-output tables and identify $\Omega_{L}$
and $\{\Omega_{s}\}_{s=1}^{S}$. The regional distribution of downstream sales, $\pi_{r}$, is
recovered from the VAT data and identifies the downstream productivity parameters $\{A_{r}\}$,
up to a multiplicative constant and therefore normalized relative to the largest commuting
zone. The extensive-margin distance coefficient of Equation~(1) identifies the trade-cost
elasticity $\alpha$ through \eqref{eq:cloglog}.

The remaining two families are where the granular structure enters. \emph{Sourcing shares} are
now aggregated to the attraction area,
\begin{equation}
\tilde\gamma_{as}=\sum_{r'\in\mathcal{R}_{a}}\tilde\gamma_{r's},
\qquad
\tilde\gamma_{r's}=\sum_{r}\frac{X_{rs}}{\sum_{r}X_{rs}}\,\gamma_{r'rs},
\label{eq:gamma-area}
\end{equation}
where $X_{rs}$ denotes the value of intermediate inputs from sector $s$ shipped to downstream
firms located in region $r$, and identify the area-level comparative advantage $\{T_{as}\}$.
For each sector these are normalized relative to the area with the largest sourcing share.
\emph{Count moments} discipline the variety count: writing $\mathcal{R}^{+}_{s}$ for the set of
regions lying inside attraction areas that host at least one supplier of sector $s$, the share
of those regions hosting no supplier is
\begin{equation}
\bar G_{s}(0)=\frac{1}{|\mathcal{R}^{+}_{s}|}\,
\#\big\{r'\in\mathcal{R}^{+}_{s}:K_{r's}=0\big\},
\qquad
\mathbb{E}\big[\bar G_{s}(0)\big]
=\frac{1}{|\mathcal{R}^{+}_{s}|}\sum_{r'\in\mathcal{R}^{+}_{s}}(1-q_{r's})^{N_{s}} .
\label{eq:count-moment}
\end{equation}
This is the count analogue of the number-of-buyers-per-seller moment in \citet{lmm2023}. It is
decreasing in $N_{s}$, which is what identifies the variety count: given the sourcing shares
--- pinned by the area aggregates and by $\alpha$, neither of which involves $N_{s}$ --- the
only free object shaping \eqref{eq:count-moment} is the number of varieties. The support is
restricted to active areas because regions in areas with no supplier at all are structural
zeros, absorbed by the fixed effect of Equation~(1) and carrying no information about
$(\alpha,N_{s})$; Appendix~\ref{app:model-zeros} makes the distinction precise.

Two observations discipline the exercise further. The variety count is bounded by the observed
number of distinct suppliers: since each variety generates one supplying firm per winning
origin, $N^{\mathrm{obs}}_{s}/R^{D}\le N_{s}\le N^{\mathrm{obs}}_{s}$ where $R^{D}$ is the
number of downstream buyers, and matching the observed firm count gives a second, independent
estimate that must fall inside those bounds (Appendix~\ref{app:model-bounds}). And because
$T_{as}$ is constant within an area, the area aggregate \eqref{eq:gamma-area} no longer
exhausts the region-level shares: each active region beyond the first in an area becomes an
over-identifying restriction, requiring that the split of an area's sourcing across its regions
follow the distance profile \eqref{eq:within-area}.

\paragraph{Identification.} The channels are close to orthogonal by construction. $\alpha$ is
read from a \emph{within}-area gradient that comparative advantage cannot touch, because
comparative advantage is constant within the area. $T_{as}$ is read from \emph{across}-area
value shares. $N_{s}$ is read from the \emph{level} of the empty-region share, which the fixed
effect strips out of the $\alpha$ regression. And $\theta$, which we calibrate, is over-identified
by the own-productivity coefficient in \eqref{eq:cloglog}.
 from the main file. Requires amsmath, amssymb and the
% \citep/\citet macros already used in the paper. Cross-references assume the
% appendix file model_appendix.tex is also included (labels app:model-*).
%
% What changed relative to the previous version of this section:
%   (i)  upstream sectors now contain a FINITE number N_s of varieties, so the
%        sourcing extensive margin is a genuine random object and observed
%        zeros are model predictions rather than dropped cells;
%   (ii) comparative advantage is located at the ATTRACTION-AREA level, the
%        level at which the extensive-margin regression identifies the
%        trade-cost elasticity;
%  (iii) trade costs are the power law tau = d^alpha, so the extensive-margin
%        distance coefficient maps into alpha with no attenuation correction.
% =====================================================================

\section{Structural model}
\label{sec:model}

We next develop a structural model that rationalizes the spatial propagation patterns
documented in the data and allows us to quantify the role of production networks in shaping
the incidence of sectoral shocks. The model has a Ricardian structure that is used to recover
firm-to-firm linkages not directly observed in the data. Based on the predicted production
network, we then simulate shocks to the downstream industry and quantify their spatial
diffusion.

Relative to a textbook Eaton--Kortum economy the model departs on one dimension: upstream
sectors contain a \emph{finite} number of varieties. This is what makes the sourcing extensive
margin --- the object our survey data measure --- a genuine random outcome. With a continuum of
varieties every region with a positive comparative advantage supplies every buyer, so the
zeros that dominate our data could only be accommodated by discarding the corresponding cells.
Here they are model predictions.

\subsection{Model}
\label{sec:model-model}

\paragraph{Environment.} The economy consists of $R+1$ regions, where the $(R+1)$th region
represents the rest of the world. Firms operate either in upstream sectors $s=1,\dots,S$ or in
the downstream industry $d$ (motor vehicles or aerospace). Each domestic region hosts at most
one representative downstream firm. Firm locations and sectoral assignments are taken as
exogenous. We allow for imports, but in the absence of balance-sheet information on foreign
suppliers we do not model firm-level linkages abroad; the foreign share of sourcing is
calibrated (Appendix~\ref{app:model-foreign}).

Domestic regions are additionally grouped into \emph{attraction areas}. An attraction area $a$
is the set of upstream regions sharing the same nearest downstream producer, and we write
$a(r')$ for the area of region $r'$ and $\mathcal{R}_{a}$ for its members. This is the same
partition as the fixed effect of the extensive-margin regression in Equation~(1), and that is
not a coincidence: it is the level at which the data identify the trade-cost elasticity.

Because the model is estimated separately for the motor vehicles and aerospace industries, we
omit the downstream-industry index in the remainder of the exposition. Unless otherwise stated,
all variables and parameters should therefore be understood as industry-specific.

\paragraph{Final demand.} Downstream firms compete under monopolistic competition and face an
isoelastic demand function:
\begin{equation}
q_{r}=\left(\frac{p_{r}}{P}\right)^{\varepsilon-1}\frac{E}{P}\,\delta^{D}_{r},
\end{equation}
where $\delta^{D}_{r}$ is a region-specific demand shifter, $E/P$ denotes aggregate real
expenditure and $p_{r}/P$ is the relative price of the variety produced in region $r$. Demand
shifters are later used to simulate downstream shocks and trace their propagation through
production networks. The market for final goods is assumed to be perfectly integrated.

\paragraph{Production.} Downstream firms produce using labor and intermediate inputs according
to a constant-returns-to-scale production technology:
\begin{equation}
y_{r}=A_{r}\,g\big(\ell_{r},\{q_{rs}\}_{s=1}^{S}\big),
\label{eq:downstream-tech}
\end{equation}
where $A_{r}$ denotes total factor productivity, $\ell_{r}$ is labor input and $q_{rs}$ is the
quantity of the input bundle from sector $s$ used in production. As discussed in
Appendix~\ref{app:model-derivation}, the structural estimation assumes a nested CES
specification, where $\Omega_{L}$ scales the labor share, $\Omega_{s}$ denotes the expenditure
share on inputs from sector $s$, and each $q_{rs}$ is itself a CES aggregate over the varieties
of sector $s$.

\paragraph{Upstream technology.} Each upstream sector $s$ consists of $N_{s}$ horizontally
differentiated varieties, indexed $\rho=1,\dots,N_{s}$. Within each variety and each region
there is a continuum of firms whose productivities above $z$ form a Poisson process with mean
$T_{r's}z^{-\theta}$. The most productive firm of region $r'$ for variety $\rho$ --- its
\emph{champion} --- therefore satisfies
\begin{equation}
\Pr\big(\varphi_{r'\rho s}\le z\big)=e^{-T_{r's}z^{-\theta}},
\qquad T_{r's}>0,
\label{eq:frechet}
\end{equation}
so one Fr\'echet draw per (region, variety) pair \emph{is} the champion of a continuum of
firms rather than an approximation to it. Because the iceberg cost scales all firms of a region
identically, that champion is the only firm of region $r'$ that ever serves a buyer sourcing
$\rho$ from $r'$: a firm in the model is a (region, variety) champion. As in Eaton and Kortum
(2002), $T_{r's}$ governs average productivity while $\theta$ controls the dispersion of
productivity draws and hence the strength of comparative advantage. Only the number of
\emph{varieties} is finite; the firms are not discretized.

\paragraph{Comparative advantage at the area level.} Comparative advantage is a property of the
attraction area rather than of the individual region:
\begin{equation}
T_{r's}=T_{a(r')s}\qquad\text{for every }r'\in\mathcal{R}_{a}.
\label{eq:T-area}
\end{equation}
Productivity draws remain independent across regions: two regions in the same area share the
Fr\'echet \emph{scale}, not the realization of their champions, so each still runs its own
Ricardian competition for every variety. What \eqref{eq:T-area} removes is unobserved
technological heterogeneity \emph{within} an area, and the payoff is that within an area the
only thing separating a region that hosts a supplier from one that does not is distance. This
is what makes the model explain the same variation the regression uses.

\paragraph{Sourcing.} Shipping intermediate goods from region $r'$ to region $r$ involves
iceberg trade costs $\tau_{r'rs}$, which we parameterize as a power law in geographic distance,
\begin{equation}
\tau_{r'rs}=d_{r'r}^{\,\alpha}.
\label{eq:tau}
\end{equation}
Following Eaton and Kortum (2002), downstream firms source each variety from the lowest-cost
supplier, in France or abroad. The probability that a downstream firm located in region $r$
sources a given variety of sector $s$ from region $r'$ is
\begin{equation}
\gamma_{r'rs}=\frac{T_{a(r')s}\,(w_{r's}\tau_{r'rs})^{-\theta}}
{\sum_{k=1}^{R+1}T_{a(k)s}\,(w_{ks}\tau_{krs})^{-\theta}}
\;\equiv\;\frac{T_{a(r')s}\,(w_{r's}\tau_{r'rs})^{-\theta}}{\Phi_{rs}},
\label{eq:ek-share}
\end{equation}
where $w_{r's}$ denotes the unit production cost in region--sector pair $(r',s)$. Because trade
costs increase with geographic distance, sourcing probabilities decline with spatial
separation, generating localized production networks. Two regions $r',r''$ in the same area
facing the same buyer satisfy
\begin{equation}
\ln\frac{\gamma_{r'rs}}{\gamma_{r''rs}}
=-\theta\alpha\big(\ln d_{r'r}-\ln d_{r''r}\big)-\theta\ln\frac{w_{r's}}{w_{r''s}},
\label{eq:within-area}
\end{equation}
the technology term and the price index having cancelled: the within-area sourcing profile is a
function of $(\theta,\alpha)$ and observed distances alone.

\paragraph{The extensive margin.} Our surveys record whether an upstream firm supplies the
downstream industry, not which buyer it serves. The relevant object is therefore not the
bilateral sourcing probability \eqref{eq:ek-share} but the probability that a region wins a
given variety \emph{somewhere} in the industry --- the union over downstream buyers,
\begin{equation}
q_{r's}\;\equiv\;\Pr\Big(\bigcup_{r=1}^{R}\big\{\text{$r'$ wins variety $\rho$ for buyer $r$}\big\}\Big).
\label{eq:q}
\end{equation}
With $N_{s}$ varieties the win events are independent across $\rho$,\footnote{Independence
across varieties follows from the granularity assumption alone (Appendix~\ref{app:model-derivation}):
the champion draws
$\varphi_{r'\rho s}$ are independent across $\rho$, and which origin wins a variety is a
comparison of \emph{relative} delivered costs, so the winner of one variety carries no
information about the winner of another. Note what this does \emph{not} require. It is silent
on the price index, because the identity of the winner is settled before any aggregation: the
sourcing decision is a per-variety cost comparison, whereas $p_{rs}$ is an aggregator built
once the winners are known. What the price index does affect is the \emph{value} attached to
each win, since expenditure shares depend on $p_{rs}(\rho)/p_{rs}$ and the realized index over
$N_{s}$ varieties is itself random. There we follow \citet{lmm2023} and let downstream firms
optimize against the expected (certainty-equivalent) index rather than internalize how the
selection of any one supplier moves it, which keeps the success probability common and
non-random across varieties in the value block as well. Appendix~\ref{app:model-price}
quantifies what this assumption costs and why it does not touch the identification of $\alpha$
or $N_{s}$.} so the industry-wide supplier count is binomial,
\begin{equation}
K_{r's}\sim\mathrm{Bin}\big(N_{s},q_{r's}\big),
\qquad
\Pr\big(K_{r's}=0\big)=(1-q_{r's})^{N_{s}}>0 .
\label{eq:binomial}
\end{equation}
The zero-supplier probability is strictly positive and endogenous: it falls with comparative
advantage, rises with distance through \eqref{eq:tau}, and falls with the variety count. Under
a continuum it is zero for every region with $T_{r's}>0$, which is the sense in which the
extensive margin previously had no content.

The union in \eqref{eq:q} is not the product of its per-buyer marginals. For a given variety
the champion cost is common across buyers and only $\tau_{r'rs}$ differs, so win events are
\emph{positively} dependent across buyers and $q_{r's}>1-\prod_{r}(1-\gamma_{r'rs})$: a firm
productive enough to win one market is more likely to win the others. The per-buyer marginal
$K_{r'rs}\sim\mathrm{Bin}(N_{s},\gamma_{r'rs})$ remains exact, but $q_{r's}$ has no tractable
closed form, and we evaluate it by simulation. This is why the model is estimated by simulated
rather than analytical method of moments.\footnote{Writing $q_{r's}$ in closed form requires
inclusion--exclusion over subsets of downstream buyers,
$q_{r's}=\sum_{\emptyset\neq\mathcal{S}\subseteq\{1,\dots,R\}}(-1)^{|\mathcal{S}|+1}
\Pr(\cap_{r\in\mathcal{S}}\{\text{$r'$ wins for }r\})$, a sum with $2^{R}-1$ terms. With the
number of downstream regions we observe this is prohibitive, and dominance pruning of the
subsets --- discarding buyers whose win event contains another's --- reduces the count but not
reliably enough to be usable inside an optimizer. The alternative of a closed form is to impose
independence and use $1-\prod_{r}(1-\gamma_{r'rs})$, which is a strict lower bound on
\eqref{eq:q} and biases the extensive margin; simulation evaluates the union exactly, at the
cost of simulation noise that the weighting matrix accounts for.}

\paragraph{Intuition: what the model implies for the extensive-margin regression.} The union
\eqref{eq:q} has no closed form, but the single-buyer case does, and it is what gives the
empirical specification of Equation~(1) its structural reading. Consider a firm of productivity
$z$ located in $r'$ facing only its area's anchor buyer $a$. Its probability of \emph{not}
supplying is complementary log-log,
\begin{equation}
\Pr\big(\text{no supply}\mid z\big)=1-\exp\big(-\exp(\eta)\big),
\qquad
\eta=\underbrace{\ln\Phi^{-r'}_{as}}_{\text{area}\times\text{sector FE}}
+\,\theta\ln w_{r's}
+\,\theta\alpha\,\ln d_{r'a}
\;-\;\theta\,\ln z ,
\label{eq:cloglog}
\end{equation}
where $\Phi^{-r'}_{as}$ collects the competitors' cost scale (Appendix~\ref{app:model-prop}
derives \eqref{eq:cloglog} and states the conditions under which it is exact). Three features
carry over to the estimation. The distance coefficient is $+\theta\alpha$, which is what makes
the extensive-margin regression informative about the trade-cost elasticity. The variety count
is \emph{absent}: conditioning on the firm removes it, so $\alpha$ and $N_{s}$ load on separate
features of the data. And the coefficient on own productivity is $-\theta$, a free
over-identifying test of a parameter we otherwise calibrate.

With several buyers the outcome is the union and \eqref{eq:cloglog} is no longer exact. We
therefore do not read $\alpha$ off the empirical coefficient directly. Instead we treat the
regression as an auxiliary model: we run the \emph{same} cloglog on simulated and observed
data, where the simulated outcome is the union event \eqref{eq:q}, and match the coefficients.
This is consistent for $\alpha$ whether or not the link is exact, and it is why the distance
coefficient enters the moment vector rather than being inverted analytically.

\subsection{Estimation and fit}
\label{sec:model-estimation}

We estimate the model using Simulated Method of Moments. Unlike the standard Ricardian trade
literature, which relies on bilateral input--output flows, such flows are not observed in our
setting. Instead, we exploit survey information on supplier relationships together with VAT
data on firm sales to discipline the spatial structure of sourcing and demand.

The model's parameters are
\begin{equation}
\big\{\Omega_{L},\;\{\Omega_{s}\},\;\alpha,\;\{A_{r}\}\big\}
\;\cup\;\big\{\{T_{as}\},\;\{N_{s}\}\big\},
\end{equation}
but only the first group is searched over. Comparative advantage and the variety count are
\emph{concentrated out}: given any candidate value of $(\Omega_{L},\{\Omega_{s}\},\alpha,\{A_{r}\})$,
each is pinned by the model itself rather than left free.

This is not a computational shortcut but a property of the moment structure. For $\{T_{as}\}$,
the area-aggregated sourcing shares \eqref{eq:gamma-area} are as numerous as the comparative-advantage
parameters, and for given trade costs and expenditures the mapping from $T$ to those shares is
invertible: there is a unique vector of comparative advantages, up to the per-sector reference
normalization, that reproduces the observed shares exactly. We compute it by a matrix-scaling
iteration alternated with the general-equilibrium solve
(Appendix~\ref{app:model-estimation}). For $\{N_{s}\}$, the count moment $\bar G_{s}(0)$ is
decreasing in the variety count and, by the argument below, does not require re-simulation, so
the variety count that best matches it is found by a monotone integer search over the bounds
\eqref{eq:app-bounds}. Both are exact concentrations in the sense of a profiled extremum
estimator: minimizing the criterion jointly over all parameters and minimizing it over
$(\Omega_{L},\{\Omega_{s}\},\alpha,\{A_{r}\})$ with $T$ and $N$ set to their concentrated values
give the same estimate.

Two things follow. The dimension of the outer search collapses from several hundred parameters
to a handful, which is what makes the estimation feasible. And the division of labour in
identification becomes explicit: the sourcing-share block is exhausted by $T$ and the count
block by $N$, so the extensive-margin coefficient is fit by $\alpha$ alone. The over-identifying
content of the area-level restriction is preserved, since a single $T_{as}$ must rationalize the
shares of every region in the area.

Beyond these, the model includes several demand- and supply-side elasticities, which we
calibrate as described in Appendix~\ref{app:model-calibration}. We also provide there
additional details on the optimization procedure and on the calculation of standard errors.

\paragraph{Targeted moments.} The model is matched to six families of moments. The first three
are as before. The downstream industry's labor share and the sectoral composition of
intermediate expenditures are constructed from input-output tables and identify $\Omega_{L}$
and $\{\Omega_{s}\}_{s=1}^{S}$. The regional distribution of downstream sales, $\pi_{r}$, is
recovered from the VAT data and identifies the downstream productivity parameters $\{A_{r}\}$,
up to a multiplicative constant and therefore normalized relative to the largest commuting
zone. The extensive-margin distance coefficient of Equation~(1) identifies the trade-cost
elasticity $\alpha$ through \eqref{eq:cloglog}.

The remaining two families are where the granular structure enters. \emph{Sourcing shares} are
now aggregated to the attraction area,
\begin{equation}
\tilde\gamma_{as}=\sum_{r'\in\mathcal{R}_{a}}\tilde\gamma_{r's},
\qquad
\tilde\gamma_{r's}=\sum_{r}\frac{X_{rs}}{\sum_{r}X_{rs}}\,\gamma_{r'rs},
\label{eq:gamma-area}
\end{equation}
where $X_{rs}$ denotes the value of intermediate inputs from sector $s$ shipped to downstream
firms located in region $r$, and identify the area-level comparative advantage $\{T_{as}\}$.
For each sector these are normalized relative to the area with the largest sourcing share.
\emph{Count moments} discipline the variety count: writing $\mathcal{R}^{+}_{s}$ for the set of
regions lying inside attraction areas that host at least one supplier of sector $s$, the share
of those regions hosting no supplier is
\begin{equation}
\bar G_{s}(0)=\frac{1}{|\mathcal{R}^{+}_{s}|}\,
\#\big\{r'\in\mathcal{R}^{+}_{s}:K_{r's}=0\big\},
\qquad
\mathbb{E}\big[\bar G_{s}(0)\big]
=\frac{1}{|\mathcal{R}^{+}_{s}|}\sum_{r'\in\mathcal{R}^{+}_{s}}(1-q_{r's})^{N_{s}} .
\label{eq:count-moment}
\end{equation}
This is the count analogue of the number-of-buyers-per-seller moment in \citet{lmm2023}. It is
decreasing in $N_{s}$, which is what identifies the variety count: given the sourcing shares
--- pinned by the area aggregates and by $\alpha$, neither of which involves $N_{s}$ --- the
only free object shaping \eqref{eq:count-moment} is the number of varieties. The support is
restricted to active areas because regions in areas with no supplier at all are structural
zeros, absorbed by the fixed effect of Equation~(1) and carrying no information about
$(\alpha,N_{s})$; Appendix~\ref{app:model-zeros} makes the distinction precise.

Two observations discipline the exercise further. The variety count is bounded by the observed
number of distinct suppliers: since each variety generates one supplying firm per winning
origin, $N^{\mathrm{obs}}_{s}/R^{D}\le N_{s}\le N^{\mathrm{obs}}_{s}$ where $R^{D}$ is the
number of downstream buyers, and matching the observed firm count gives a second, independent
estimate that must fall inside those bounds (Appendix~\ref{app:model-bounds}). And because
$T_{as}$ is constant within an area, the area aggregate \eqref{eq:gamma-area} no longer
exhausts the region-level shares: each active region beyond the first in an area becomes an
over-identifying restriction, requiring that the split of an area's sourcing across its regions
follow the distance profile \eqref{eq:within-area}.

\paragraph{Identification.} The channels are close to orthogonal by construction. $\alpha$ is
read from a \emph{within}-area gradient that comparative advantage cannot touch, because
comparative advantage is constant within the area. $T_{as}$ is read from \emph{across}-area
value shares. $N_{s}$ is read from the \emph{level} of the empty-region share, which the fixed
effect strips out of the $\alpha$ regression. And $\theta$, which we calibrate, is over-identified
by the own-productivity coefficient in \eqref{eq:cloglog}.
 from the main file. Requires amsmath, amssymb and the
% \citep/\citet macros already used in the paper. Cross-references assume the
% appendix file model_appendix.tex is also included (labels app:model-*).
%
% What changed relative to the previous version of this section:
%   (i)  upstream sectors now contain a FINITE number N_s of varieties, so the
%        sourcing extensive margin is a genuine random object and observed
%        zeros are model predictions rather than dropped cells;
%   (ii) comparative advantage is located at the ATTRACTION-AREA level, the
%        level at which the extensive-margin regression identifies the
%        trade-cost elasticity;
%  (iii) trade costs are the power law tau = d^alpha, so the extensive-margin
%        distance coefficient maps into alpha with no attenuation correction.
% =====================================================================

\section{Structural model}
\label{sec:model}

We next develop a structural model that rationalizes the spatial propagation patterns
documented in the data and allows us to quantify the role of production networks in shaping
the incidence of sectoral shocks. The model has a Ricardian structure that is used to recover
firm-to-firm linkages not directly observed in the data. Based on the predicted production
network, we then simulate shocks to the downstream industry and quantify their spatial
diffusion.

Relative to a textbook Eaton--Kortum economy the model departs on one dimension: upstream
sectors contain a \emph{finite} number of varieties. This is what makes the sourcing extensive
margin --- the object our survey data measure --- a genuine random outcome. With a continuum of
varieties every region with a positive comparative advantage supplies every buyer, so the
zeros that dominate our data could only be accommodated by discarding the corresponding cells.
Here they are model predictions.

\subsection{Model}
\label{sec:model-model}

\paragraph{Environment.} The economy consists of $R+1$ regions, where the $(R+1)$th region
represents the rest of the world. Firms operate either in upstream sectors $s=1,\dots,S$ or in
the downstream industry $d$ (motor vehicles or aerospace). Each domestic region hosts at most
one representative downstream firm. Firm locations and sectoral assignments are taken as
exogenous. We allow for imports, but in the absence of balance-sheet information on foreign
suppliers we do not model firm-level linkages abroad; the foreign share of sourcing is
calibrated (Appendix~\ref{app:model-foreign}).

Domestic regions are additionally grouped into \emph{attraction areas}. An attraction area $a$
is the set of upstream regions sharing the same nearest downstream producer, and we write
$a(r')$ for the area of region $r'$ and $\mathcal{R}_{a}$ for its members. This is the same
partition as the fixed effect of the extensive-margin regression in Equation~(1), and that is
not a coincidence: it is the level at which the data identify the trade-cost elasticity.

Because the model is estimated separately for the motor vehicles and aerospace industries, we
omit the downstream-industry index in the remainder of the exposition. Unless otherwise stated,
all variables and parameters should therefore be understood as industry-specific.

\paragraph{Final demand.} Downstream firms compete under monopolistic competition and face an
isoelastic demand function:
\begin{equation}
q_{r}=\left(\frac{p_{r}}{P}\right)^{\varepsilon-1}\frac{E}{P}\,\delta^{D}_{r},
\end{equation}
where $\delta^{D}_{r}$ is a region-specific demand shifter, $E/P$ denotes aggregate real
expenditure and $p_{r}/P$ is the relative price of the variety produced in region $r$. Demand
shifters are later used to simulate downstream shocks and trace their propagation through
production networks. The market for final goods is assumed to be perfectly integrated.

\paragraph{Production.} Downstream firms produce using labor and intermediate inputs according
to a constant-returns-to-scale production technology:
\begin{equation}
y_{r}=A_{r}\,g\big(\ell_{r},\{q_{rs}\}_{s=1}^{S}\big),
\label{eq:downstream-tech}
\end{equation}
where $A_{r}$ denotes total factor productivity, $\ell_{r}$ is labor input and $q_{rs}$ is the
quantity of the input bundle from sector $s$ used in production. As discussed in
Appendix~\ref{app:model-derivation}, the structural estimation assumes a nested CES
specification, where $\Omega_{L}$ scales the labor share, $\Omega_{s}$ denotes the expenditure
share on inputs from sector $s$, and each $q_{rs}$ is itself a CES aggregate over the varieties
of sector $s$.

\paragraph{Upstream technology.} Each upstream sector $s$ consists of $N_{s}$ horizontally
differentiated varieties, indexed $\rho=1,\dots,N_{s}$. Within each variety and each region
there is a continuum of firms whose productivities above $z$ form a Poisson process with mean
$T_{r's}z^{-\theta}$. The most productive firm of region $r'$ for variety $\rho$ --- its
\emph{champion} --- therefore satisfies
\begin{equation}
\Pr\big(\varphi_{r'\rho s}\le z\big)=e^{-T_{r's}z^{-\theta}},
\qquad T_{r's}>0,
\label{eq:frechet}
\end{equation}
so one Fr\'echet draw per (region, variety) pair \emph{is} the champion of a continuum of
firms rather than an approximation to it. Because the iceberg cost scales all firms of a region
identically, that champion is the only firm of region $r'$ that ever serves a buyer sourcing
$\rho$ from $r'$: a firm in the model is a (region, variety) champion. As in Eaton and Kortum
(2002), $T_{r's}$ governs average productivity while $\theta$ controls the dispersion of
productivity draws and hence the strength of comparative advantage. Only the number of
\emph{varieties} is finite; the firms are not discretized.

\paragraph{Comparative advantage at the area level.} Comparative advantage is a property of the
attraction area rather than of the individual region:
\begin{equation}
T_{r's}=T_{a(r')s}\qquad\text{for every }r'\in\mathcal{R}_{a}.
\label{eq:T-area}
\end{equation}
Productivity draws remain independent across regions: two regions in the same area share the
Fr\'echet \emph{scale}, not the realization of their champions, so each still runs its own
Ricardian competition for every variety. What \eqref{eq:T-area} removes is unobserved
technological heterogeneity \emph{within} an area, and the payoff is that within an area the
only thing separating a region that hosts a supplier from one that does not is distance. This
is what makes the model explain the same variation the regression uses.

\paragraph{Sourcing.} Shipping intermediate goods from region $r'$ to region $r$ involves
iceberg trade costs $\tau_{r'rs}$, which we parameterize as a power law in geographic distance,
\begin{equation}
\tau_{r'rs}=d_{r'r}^{\,\alpha}.
\label{eq:tau}
\end{equation}
Following Eaton and Kortum (2002), downstream firms source each variety from the lowest-cost
supplier, in France or abroad. The probability that a downstream firm located in region $r$
sources a given variety of sector $s$ from region $r'$ is
\begin{equation}
\gamma_{r'rs}=\frac{T_{a(r')s}\,(w_{r's}\tau_{r'rs})^{-\theta}}
{\sum_{k=1}^{R+1}T_{a(k)s}\,(w_{ks}\tau_{krs})^{-\theta}}
\;\equiv\;\frac{T_{a(r')s}\,(w_{r's}\tau_{r'rs})^{-\theta}}{\Phi_{rs}},
\label{eq:ek-share}
\end{equation}
where $w_{r's}$ denotes the unit production cost in region--sector pair $(r',s)$. Because trade
costs increase with geographic distance, sourcing probabilities decline with spatial
separation, generating localized production networks. Two regions $r',r''$ in the same area
facing the same buyer satisfy
\begin{equation}
\ln\frac{\gamma_{r'rs}}{\gamma_{r''rs}}
=-\theta\alpha\big(\ln d_{r'r}-\ln d_{r''r}\big)-\theta\ln\frac{w_{r's}}{w_{r''s}},
\label{eq:within-area}
\end{equation}
the technology term and the price index having cancelled: the within-area sourcing profile is a
function of $(\theta,\alpha)$ and observed distances alone.

\paragraph{The extensive margin.} Our surveys record whether an upstream firm supplies the
downstream industry, not which buyer it serves. The relevant object is therefore not the
bilateral sourcing probability \eqref{eq:ek-share} but the probability that a region wins a
given variety \emph{somewhere} in the industry --- the union over downstream buyers,
\begin{equation}
q_{r's}\;\equiv\;\Pr\Big(\bigcup_{r=1}^{R}\big\{\text{$r'$ wins variety $\rho$ for buyer $r$}\big\}\Big).
\label{eq:q}
\end{equation}
With $N_{s}$ varieties the win events are independent across $\rho$,\footnote{Independence
across varieties follows from the granularity assumption alone (Appendix~\ref{app:model-derivation}):
the champion draws
$\varphi_{r'\rho s}$ are independent across $\rho$, and which origin wins a variety is a
comparison of \emph{relative} delivered costs, so the winner of one variety carries no
information about the winner of another. Note what this does \emph{not} require. It is silent
on the price index, because the identity of the winner is settled before any aggregation: the
sourcing decision is a per-variety cost comparison, whereas $p_{rs}$ is an aggregator built
once the winners are known. What the price index does affect is the \emph{value} attached to
each win, since expenditure shares depend on $p_{rs}(\rho)/p_{rs}$ and the realized index over
$N_{s}$ varieties is itself random. There we follow \citet{lmm2023} and let downstream firms
optimize against the expected (certainty-equivalent) index rather than internalize how the
selection of any one supplier moves it, which keeps the success probability common and
non-random across varieties in the value block as well. Appendix~\ref{app:model-price}
quantifies what this assumption costs and why it does not touch the identification of $\alpha$
or $N_{s}$.} so the industry-wide supplier count is binomial,
\begin{equation}
K_{r's}\sim\mathrm{Bin}\big(N_{s},q_{r's}\big),
\qquad
\Pr\big(K_{r's}=0\big)=(1-q_{r's})^{N_{s}}>0 .
\label{eq:binomial}
\end{equation}
The zero-supplier probability is strictly positive and endogenous: it falls with comparative
advantage, rises with distance through \eqref{eq:tau}, and falls with the variety count. Under
a continuum it is zero for every region with $T_{r's}>0$, which is the sense in which the
extensive margin previously had no content.

The union in \eqref{eq:q} is not the product of its per-buyer marginals. For a given variety
the champion cost is common across buyers and only $\tau_{r'rs}$ differs, so win events are
\emph{positively} dependent across buyers and $q_{r's}>1-\prod_{r}(1-\gamma_{r'rs})$: a firm
productive enough to win one market is more likely to win the others. The per-buyer marginal
$K_{r'rs}\sim\mathrm{Bin}(N_{s},\gamma_{r'rs})$ remains exact, but $q_{r's}$ has no tractable
closed form, and we evaluate it by simulation. This is why the model is estimated by simulated
rather than analytical method of moments.\footnote{Writing $q_{r's}$ in closed form requires
inclusion--exclusion over subsets of downstream buyers,
$q_{r's}=\sum_{\emptyset\neq\mathcal{S}\subseteq\{1,\dots,R\}}(-1)^{|\mathcal{S}|+1}
\Pr(\cap_{r\in\mathcal{S}}\{\text{$r'$ wins for }r\})$, a sum with $2^{R}-1$ terms. With the
number of downstream regions we observe this is prohibitive, and dominance pruning of the
subsets --- discarding buyers whose win event contains another's --- reduces the count but not
reliably enough to be usable inside an optimizer. The alternative of a closed form is to impose
independence and use $1-\prod_{r}(1-\gamma_{r'rs})$, which is a strict lower bound on
\eqref{eq:q} and biases the extensive margin; simulation evaluates the union exactly, at the
cost of simulation noise that the weighting matrix accounts for.}

\paragraph{Intuition: what the model implies for the extensive-margin regression.} The union
\eqref{eq:q} has no closed form, but the single-buyer case does, and it is what gives the
empirical specification of Equation~(1) its structural reading. Consider a firm of productivity
$z$ located in $r'$ facing only its area's anchor buyer $a$. Its probability of \emph{not}
supplying is complementary log-log,
\begin{equation}
\Pr\big(\text{no supply}\mid z\big)=1-\exp\big(-\exp(\eta)\big),
\qquad
\eta=\underbrace{\ln\Phi^{-r'}_{as}}_{\text{area}\times\text{sector FE}}
+\,\theta\ln w_{r's}
+\,\theta\alpha\,\ln d_{r'a}
\;-\;\theta\,\ln z ,
\label{eq:cloglog}
\end{equation}
where $\Phi^{-r'}_{as}$ collects the competitors' cost scale (Appendix~\ref{app:model-prop}
derives \eqref{eq:cloglog} and states the conditions under which it is exact). Three features
carry over to the estimation. The distance coefficient is $+\theta\alpha$, which is what makes
the extensive-margin regression informative about the trade-cost elasticity. The variety count
is \emph{absent}: conditioning on the firm removes it, so $\alpha$ and $N_{s}$ load on separate
features of the data. And the coefficient on own productivity is $-\theta$, a free
over-identifying test of a parameter we otherwise calibrate.

With several buyers the outcome is the union and \eqref{eq:cloglog} is no longer exact. We
therefore do not read $\alpha$ off the empirical coefficient directly. Instead we treat the
regression as an auxiliary model: we run the \emph{same} cloglog on simulated and observed
data, where the simulated outcome is the union event \eqref{eq:q}, and match the coefficients.
This is consistent for $\alpha$ whether or not the link is exact, and it is why the distance
coefficient enters the moment vector rather than being inverted analytically.

\subsection{Estimation and fit}
\label{sec:model-estimation}

We estimate the model using Simulated Method of Moments. Unlike the standard Ricardian trade
literature, which relies on bilateral input--output flows, such flows are not observed in our
setting. Instead, we exploit survey information on supplier relationships together with VAT
data on firm sales to discipline the spatial structure of sourcing and demand.

The model's parameters are
\begin{equation}
\big\{\Omega_{L},\;\{\Omega_{s}\},\;\alpha,\;\{A_{r}\}\big\}
\;\cup\;\big\{\{T_{as}\},\;\{N_{s}\}\big\},
\end{equation}
but only the first group is searched over. Comparative advantage and the variety count are
\emph{concentrated out}: given any candidate value of $(\Omega_{L},\{\Omega_{s}\},\alpha,\{A_{r}\})$,
each is pinned by the model itself rather than left free.

This is not a computational shortcut but a property of the moment structure. For $\{T_{as}\}$,
the area-aggregated sourcing shares \eqref{eq:gamma-area} are as numerous as the comparative-advantage
parameters, and for given trade costs and expenditures the mapping from $T$ to those shares is
invertible: there is a unique vector of comparative advantages, up to the per-sector reference
normalization, that reproduces the observed shares exactly. We compute it by a matrix-scaling
iteration alternated with the general-equilibrium solve
(Appendix~\ref{app:model-estimation}). For $\{N_{s}\}$, the count moment $\bar G_{s}(0)$ is
decreasing in the variety count and, by the argument below, does not require re-simulation, so
the variety count that best matches it is found by a monotone integer search over the bounds
\eqref{eq:app-bounds}. Both are exact concentrations in the sense of a profiled extremum
estimator: minimizing the criterion jointly over all parameters and minimizing it over
$(\Omega_{L},\{\Omega_{s}\},\alpha,\{A_{r}\})$ with $T$ and $N$ set to their concentrated values
give the same estimate.

Two things follow. The dimension of the outer search collapses from several hundred parameters
to a handful, which is what makes the estimation feasible. And the division of labour in
identification becomes explicit: the sourcing-share block is exhausted by $T$ and the count
block by $N$, so the extensive-margin coefficient is fit by $\alpha$ alone. The over-identifying
content of the area-level restriction is preserved, since a single $T_{as}$ must rationalize the
shares of every region in the area.

Beyond these, the model includes several demand- and supply-side elasticities, which we
calibrate as described in Appendix~\ref{app:model-calibration}. We also provide there
additional details on the optimization procedure and on the calculation of standard errors.

\paragraph{Targeted moments.} The model is matched to six families of moments. The first three
are as before. The downstream industry's labor share and the sectoral composition of
intermediate expenditures are constructed from input-output tables and identify $\Omega_{L}$
and $\{\Omega_{s}\}_{s=1}^{S}$. The regional distribution of downstream sales, $\pi_{r}$, is
recovered from the VAT data and identifies the downstream productivity parameters $\{A_{r}\}$,
up to a multiplicative constant and therefore normalized relative to the largest commuting
zone. The extensive-margin distance coefficient of Equation~(1) identifies the trade-cost
elasticity $\alpha$ through \eqref{eq:cloglog}.

The remaining two families are where the granular structure enters. \emph{Sourcing shares} are
now aggregated to the attraction area,
\begin{equation}
\tilde\gamma_{as}=\sum_{r'\in\mathcal{R}_{a}}\tilde\gamma_{r's},
\qquad
\tilde\gamma_{r's}=\sum_{r}\frac{X_{rs}}{\sum_{r}X_{rs}}\,\gamma_{r'rs},
\label{eq:gamma-area}
\end{equation}
where $X_{rs}$ denotes the value of intermediate inputs from sector $s$ shipped to downstream
firms located in region $r$, and identify the area-level comparative advantage $\{T_{as}\}$.
For each sector these are normalized relative to the area with the largest sourcing share.
\emph{Count moments} discipline the variety count: writing $\mathcal{R}^{+}_{s}$ for the set of
regions lying inside attraction areas that host at least one supplier of sector $s$, the share
of those regions hosting no supplier is
\begin{equation}
\bar G_{s}(0)=\frac{1}{|\mathcal{R}^{+}_{s}|}\,
\#\big\{r'\in\mathcal{R}^{+}_{s}:K_{r's}=0\big\},
\qquad
\mathbb{E}\big[\bar G_{s}(0)\big]
=\frac{1}{|\mathcal{R}^{+}_{s}|}\sum_{r'\in\mathcal{R}^{+}_{s}}(1-q_{r's})^{N_{s}} .
\label{eq:count-moment}
\end{equation}
This is the count analogue of the number-of-buyers-per-seller moment in \citet{lmm2023}. It is
decreasing in $N_{s}$, which is what identifies the variety count: given the sourcing shares
--- pinned by the area aggregates and by $\alpha$, neither of which involves $N_{s}$ --- the
only free object shaping \eqref{eq:count-moment} is the number of varieties. The support is
restricted to active areas because regions in areas with no supplier at all are structural
zeros, absorbed by the fixed effect of Equation~(1) and carrying no information about
$(\alpha,N_{s})$; Appendix~\ref{app:model-zeros} makes the distinction precise.

Two observations discipline the exercise further. The variety count is bounded by the observed
number of distinct suppliers: since each variety generates one supplying firm per winning
origin, $N^{\mathrm{obs}}_{s}/R^{D}\le N_{s}\le N^{\mathrm{obs}}_{s}$ where $R^{D}$ is the
number of downstream buyers, and matching the observed firm count gives a second, independent
estimate that must fall inside those bounds (Appendix~\ref{app:model-bounds}). And because
$T_{as}$ is constant within an area, the area aggregate \eqref{eq:gamma-area} no longer
exhausts the region-level shares: each active region beyond the first in an area becomes an
over-identifying restriction, requiring that the split of an area's sourcing across its regions
follow the distance profile \eqref{eq:within-area}.

\paragraph{Identification.} The channels are close to orthogonal by construction. $\alpha$ is
read from a \emph{within}-area gradient that comparative advantage cannot touch, because
comparative advantage is constant within the area. $T_{as}$ is read from \emph{across}-area
value shares. $N_{s}$ is read from the \emph{level} of the empty-region share, which the fixed
effect strips out of the $\alpha$ regression. And $\theta$, which we calibrate, is over-identified
by the own-productivity coefficient in \eqref{eq:cloglog}.
